\chapter{Implementing a GPT Model}\label{ch:gpt}

In \cref{ch:text} we turned raw text into tensors of token IDs and embeddings, and in
\cref{ch:attention} we built the causal multi-head attention module that lets every token
look back at its predecessors. Those are the two hardest ingredients. This chapter assembles
them, together with a handful of small but essential components (layer normalization, a
position-wise feed-forward network, residual connections, dropout), into a complete
GPT-2-style language model that maps a batch of token IDs to next-token logits, and then
uses that model to generate text.

For a systems person this chapter has a second, equally important goal: \emph{accounting}.
Once the model exists as code we can say \emph{exactly} how many parameters it has, how many
floating-point operations a forward pass costs, how many bytes of activations the backward pass
will need, and which kernels dominate the runtime. Every later systems decision in this
book---how big a batch fits on a GPU, when decoding becomes memory bound, how to shard a model
across devices, how to size a KV cache---rests on the arithmetic we derive here. We will check
that arithmetic against PyTorch itself, using \code{FlopCounterMode} and autograd's
saved-tensor hooks, so that nothing is taken on faith.

\begin{figure}[t]
\centering
\resizebox{\linewidth}{!}{%
\begin{tikzpicture}[
  stage/.style={draw=sysblue,fill=sysbluebg,rounded corners=2pt,minimum width=2.3cm,
    minimum height=1.05cm,align=center,font=\small},
  sysst/.style={stage,draw=napkin,fill=napkinbg},
  num/.style={circle,fill=sysblue,text=white,font=\scriptsize\bfseries,inner sep=1.2pt},
  arr/.style={-{Stealth[length=2.2mm]},thick,sysblue!70}]
\node[stage] (s1) {Skeleton \&\\config};
\node[stage,right=0.45cm of s1] (s2) {Layer\\normalization};
\node[stage,right=0.45cm of s2] (s3) {GELU\\feed-forward};
\node[stage,right=0.45cm of s3] (s4) {Residual\\connections};
\node[stage,below=0.9cm of s4] (s5) {Transformer\\block};
\node[stage,left=0.45cm of s5] (s6) {Full GPT\\model};
\node[stage,left=0.45cm of s6] (s7) {Greedy\\generation};
\node[sysst,left=0.45cm of s7] (s8) {Accounting:\\params, FLOPs,\\bytes};
\foreach \s/\n in {s1/1,s2/2,s3/3,s4/4,s5/5,s6/6,s7/7,s8/8}
  \node[num] at (\s.north west) {\n};
\draw[arr] (s1)--(s2); \draw[arr] (s2)--(s3); \draw[arr] (s3)--(s4);
\draw[arr] (s4)--(s5); \draw[arr] (s5)--(s6); \draw[arr] (s6)--(s7); \draw[arr] (s7)--(s8);
\node[font=\scriptsize\itshape,text=srcgray,above=0.1cm of s2.north east,anchor=south west]
  {components (\cref{sec:gpt:ln,sec:gpt:ffn,sec:gpt:shortcut})};
\node[font=\scriptsize\itshape,text=napkin,below=0.12cm of s8] {systems layer (\cref{sec:gpt:accounting})};
\end{tikzpicture}}
\caption{Roadmap of this chapter. We first fix the overall skeleton and configuration, build
the three missing components, compose them into a transformer block and then a full model,
generate text with it, and finally account for every parameter, FLOP and byte.}
\label{fig:gpt:roadmap}
\end{figure}

% =====================================================================
\section{Coding an LLM architecture}\label{sec:gpt:arch}

A GPT (\emph{generative pretrained transformer}) is a decoder-only transformer: a stack of
identical blocks, each combining causal self-attention with a small per-token neural network,
sandwiched between an embedding layer at the bottom and a linear ``unembedding'' layer at the
top. The original GPT-2 release (Radford et al., 2019) came in four sizes. We will build the
smallest, usually called the 124M model, but every line of code is size-agnostic: the four
variants differ only in three numbers.

\begin{table}[h]
\centering\small
\begin{tabular}{lrrrrr}
\toprule
Variant & $\dmodel$ & Layers $L$ & Heads $H$ & Params (tied) & Params (untied, no QKV bias)\\
\midrule
GPT-2 small  & 768  & 12 & 12 & \num{124439808}  & \num{163009536}\\
GPT-2 medium & 1024 & 24 & 16 & \num{354823168}  & \num{406212608}\\
GPT-2 large  & 1280 & 36 & 20 & \num{774030080}  & \num{838220800}\\
GPT-2 XL     & 1600 & 48 & 25 & \num{1557611200} & \num{1637792000}\\
\bottomrule
\end{tabular}
\caption{The GPT-2 family. All share a vocabulary of \num{50257} tokens and a context of
1024 positions. ``Tied'' counts the released checkpoints, which share the token-embedding and
output matrices and use biases on the query/key/value projections; the last column is our
from-scratch configuration with a separate output head. Counts computed by the Lab~5 code.}
\label{tab:gpt:sizes}
\end{table}

We describe a model by a plain dictionary. Keeping the configuration as data (rather than
hard-coding numbers) is what lets the same class instantiate a 64-dimensional toy model for
unit tests and the real 1.5B-parameter model for inference.

\begin{lstlisting}[caption={The configuration that defines GPT-2 small in our code.},label={lst:gpt:config}]
GPT_CONFIG_124M = dict(
    vocab_size=50257,     # BPE vocabulary of the GPT-2 tokenizer (Ch. 3)
    context_length=1024,  # number of learned position embeddings
    emb_dim=768,          # d_model: width of the residual stream
    n_heads=12,           # attention heads per layer, head_dim = 768/12 = 64
    n_layers=12,          # number of transformer blocks
    drop_rate=0.1,        # dropout probability (0.0 is common for fine-tuning)
    qkv_bias=False,       # Q/K/V biases (True for the OpenAI weights)
)
\end{lstlisting}

\paragraph{Shapes are the contract.}
Before writing any layer it pays to fix the shape of the data at every interface, because the
whole model is just a sequence of shape-preserving transformations on a tensor of width
$\dmodel$. \Cref{fig:gpt:topdown} shows the top-down view. A batch of $B$ sequences of $T$ token
IDs, shape $(B,T)$, is embedded into $(B,T,\dmodel)$ by adding a token embedding and a learned
position embedding. That tensor---the \emph{residual stream}---flows through $L$ transformer
blocks, each of which reads it and writes an update of the same shape. A final normalization and
a linear map to the vocabulary produce logits of shape $(B,T,\vocab)$: one score per vocabulary
entry, for every position.

\begin{figure}[t]
\centering
\begin{tikzpicture}[
  blk/.style={draw,rounded corners=2pt,minimum width=5.2cm,minimum height=0.7cm,font=\small,align=center},
  shp/.style={font=\scriptsize\ttfamily,text=srcgray,anchor=west},
  arr/.style={-{Stealth[length=2mm]},thick}]
\node[font=\small] (txt) at (0,0) {``Every effort moves you''};
\node[blk,fill=fdata!15,above=0.45cm of txt] (tok) {Tokenizer (BPE, \cref{ch:text})};
\node[blk,fill=fdata!25,above=0.45cm of tok] (emb) {Token embedding \ $+$\ position embedding};
\node[blk,fill=white,above=0.35cm of emb] (drop) {Dropout};
\node[blk,fill=fcompute!20,minimum height=1.5cm,above=0.45cm of drop] (trf)
  {Transformer block\\[2pt]\footnotesize LN $\to$ masked MHA $\to$ $+$\ \ LN $\to$ FFN $\to$ $+$};
\node[blk,fill=white,above=0.45cm of trf] (fln) {Final LayerNorm};
\node[blk,fill=fcompute!35,above=0.35cm of fln] (head) {Linear output head $\dmodel\to\vocab$};
\node[font=\small,above=0.45cm of head] (out) {next-token logits for every position};
\draw[arr] (txt)--(tok); \draw[arr] (tok)--(emb); \draw[arr] (emb)--(drop);
\draw[arr] (drop)--(trf); \draw[arr] (trf)--(fln); \draw[arr] (fln)--(head); \draw[arr] (head)--(out);
\node[shp] at ($(tok.east)+(0.25,-0.32)$) {(B, T) int64};
\node[shp] at ($(emb.east)+(0.25,0)$) {(B, T, d)};
\node[shp] at ($(trf.east)+(0.25,0)$) {(B, T, d)};
\node[shp] at ($(head.east)+(0.25,0)$) {(B, T, V)};
\draw[decorate,decoration={brace,amplitude=5pt},thick,sysblue]
  ([xshift=-0.15cm]trf.south west) -- node[left=6pt,font=\small,text=sysblue]{$\times L$} ([xshift=-0.15cm]trf.north west);
\end{tikzpicture}
\caption{Top-down view of a GPT model with the tensor shape at each interface. Everything
between the embedding and the output head preserves the shape $(B,T,\dmodel)$.}
\label{fig:gpt:topdown}
\end{figure}

\paragraph{Stub-first development.}
A good way to build any multi-component system is to write the outermost structure first with
placeholder components that already honor the interfaces, run it end to end, and then replace
the placeholders one by one. In PyTorch \code{nn.Identity} is a perfect placeholder for a block
or a normalization layer: it accepts a tensor of any shape and returns it unchanged. The
skeleton in \cref{lst:gpt:skeleton} runs, produces correctly shaped logits, and gives us a
harness to test each real component as it arrives.

\begin{lstlisting}[caption={A runnable skeleton: real embeddings and head, placeholder
blocks and norm.},label={lst:gpt:skeleton}]
import torch, torch.nn as nn

class SkeletonGPT(nn.Module):
    def __init__(self, cfg):
        super().__init__()
        self.tok_emb = nn.Embedding(cfg["vocab_size"], cfg["emb_dim"])
        self.pos_emb = nn.Embedding(cfg["context_length"], cfg["emb_dim"])
        self.drop_emb = nn.Dropout(cfg["drop_rate"])
        # placeholders that respect the (B, T, d) -> (B, T, d) contract
        stubs = [nn.Identity() for _ in range(cfg["n_layers"])]
        self.trf_blocks = nn.Sequential(*stubs)
        self.final_norm = nn.Identity()
        self.out_head = nn.Linear(cfg["emb_dim"], cfg["vocab_size"], bias=False)

    def forward(self, idx):                       # idx: (B, T)
        T = idx.shape[1]
        x = self.tok_emb(idx) + self.pos_emb(torch.arange(T, device=idx.device))
        x = self.trf_blocks(self.drop_emb(x))      # (B, T, d)
        return self.out_head(self.final_norm(x))   # (B, T, V)

logits = SkeletonGPT(GPT_CONFIG_124M)(torch.randint(0, 50257, (2, 4)))
print(logits.shape)   # torch.Size([2, 4, 50257])
\end{lstlisting}

Each of the $2\times4$ output rows is a 50,257-dimensional vector. Position $t$ of row $b$ holds
the model's (currently random) scores for the token that should follow the first $t+1$ tokens of
sequence $b$. During training \emph{every} position contributes a prediction, which is why a
single sequence of length $T$ yields $T$ training examples.

\begin{systemslens}[The vocabulary is a systems parameter]
Two things about $\vocab=\num{50257}$ matter for performance. First, the output head is the single
largest matrix in GPT-2 small: $\vocab\dmodel=\num{38597376}$ parameters, 24\% of the untied model.
Second, \num{50257} is an awkward GEMM dimension: it is odd, so tiles of 64 or 128 columns on a
tensor-core GPU always leave a ragged remainder. Padding the vocabulary to \num{50304}
($=393\times128$) adds only \num{36096} never-used parameters but lets every tile be full;
nanoGPT-style training scripts do exactly this and observe a measurable speedup. The general rule:
choose $\dmodel$, $\dff$, head dimension and vocabulary as multiples of 64 (better 128) so the
matmul kernels never run partially empty tiles.
\end{systemslens}

\begin{napkinmath}[The logits tensor is bigger than the model]
For a training batch of $B=8$ sequences of $T=1024$ tokens, the logits have
$8\times1024\times\num{50257}\approx\num{4.1e8}$ entries. In fp32 that is
$\SI{1.65}{\giga\byte}$---more than twice the \SI{0.62}{\giga\byte} of fp32 weights of the
untied 163M model---and the softmax inside the cross-entropy loss materializes another tensor of
the same size for the backward pass. This is why large-vocabulary models (Qwen3 has
$\vocab\approx\num{152000}$) compute the loss in \emph{chunks} of the sequence, or with fused
``linear + cross-entropy'' kernels that never store the full logits (\cref{ch:pretrain}).
\end{napkinmath}

% =====================================================================
\section{Normalizing activations with layer normalization}\label{sec:gpt:ln}

Deep networks are hard to train when the scale of the activations drifts from layer to layer:
gradients either explode or shrink exponentially with depth, and the optimizer spends its
effort chasing scale rather than learning features. \emph{Layer normalization} (Ba et al., 2016)
fixes the scale locally. For each token vector $\vect{x}\in\R^{\dmodel}$ independently, it
subtracts the mean of the features and divides by their standard deviation, then applies a
learned per-feature affine map:
\begin{equation}
\LN(\vect{x}) = \vect{\gamma}\odot\frac{\vect{x}-\mu}{\sqrt{\sigma^2+\epsilon}}+\vect{\beta},
\qquad \mu=\frac{1}{d}\sum_{i=1}^{d}x_i,\quad \sigma^2=\frac{1}{d}\sum_{i=1}^{d}(x_i-\mu)^2 .
\label{eq:gpt:ln}
\end{equation}
The trainable \emph{scale} $\vect\gamma$ (initialized to ones) and \emph{shift} $\vect\beta$
(initialized to zeros) let the network undo the normalization where that helps; the small
$\epsilon=10^{-5}$ prevents division by zero for a constant input.

\paragraph{Along which axis?}
The statistics are computed over the \emph{feature} dimension, which in a $(B,T,\dmodel)$ tensor
is the last one, so the code uses \code{dim=-1} with \code{keepdim=True} to keep the result
broadcastable. \Cref{fig:gpt:normaxes} contrasts this with batch normalization, which reduces
over the batch (and sequence) axes instead. The difference is decisive for language models:
layer norm treats each token in isolation, so its output does not depend on what else happens to
be in the batch, on the padding, or on how many GPUs the batch is split across.

\begin{figure}[t]
\centering
\begin{tikzpicture}[x=0.42cm,y=0.42cm]
\providecommand{\gptnormgrid}{}\renewcommand{\gptnormgrid}[3]{%
  \begin{scope}[shift={(#1,0)}]
  \foreach \i in {0,...,5}{\foreach \j in {0,...,3}{
    \draw[gray!60,fill=white] (\i,\j) rectangle ++(1,1);}}
  #2
  \node[font=\scriptsize,anchor=north,align=center,text width=3.6cm] at (3,-0.5) {#3};
  \draw[-{Stealth},thick] (-0.5,0)--(-0.5,4.4) node[above,font=\scriptsize]{tokens};
  \draw[-{Stealth},thick] (0,-0.0)--(6.6,-0.0) node[right,font=\scriptsize]{features};
  \end{scope}}
\gptnormgrid{0}{\foreach \i in {0,...,5}{\fill[keygreen!45] (\i,2) rectangle ++(1,1);}
  \draw[keygreen,very thick] (0,2) rectangle (6,3);
  \node[font=\scriptsize,text=keygreen,anchor=west] at (6.2,2.5) {one $\mu,\sigma$};}
  {LayerNorm: per token, over features}
\gptnormgrid{15}{\foreach \j in {0,...,3}{\fill[warnred!35] (2,\j) rectangle ++(1,1);}
  \draw[warnred,very thick] (2,0) rectangle (3,4);
  \node[font=\scriptsize,text=warnred,anchor=west] at (3.2,2.5) {one $\mu,\sigma$};}
  {BatchNorm: per feature, over tokens/batch}
\end{tikzpicture}
\caption{Which entries share normalization statistics. Layer norm (left) computes one mean and
variance per token across its $\dmodel$ features; batch norm (right) computes one per feature
across all tokens in the batch, which couples examples together.}
\label{fig:gpt:normaxes}
\end{figure}

\begin{lstlisting}[caption={Layer normalization written out explicitly (as in \ident{llmsys/layers.py}).},label={lst:gpt:ln}]
class LayerNorm(nn.Module):
    def __init__(self, emb_dim, eps=1e-5):
        super().__init__()
        self.eps = eps
        self.scale = nn.Parameter(torch.ones(emb_dim))    # gamma
        self.shift = nn.Parameter(torch.zeros(emb_dim))   # beta

    def forward(self, x):
        mean = x.mean(dim=-1, keepdim=True)
        var = x.var(dim=-1, keepdim=True, unbiased=False)  # divide by d, not d-1
        return self.scale * (x - mean) / torch.sqrt(var + self.eps) + self.shift
\end{lstlisting}

Two details deserve comment. We use the \emph{biased} variance estimator (dividing by $d$
rather than $d-1$). Bessel's correction is meaningless here because we are not estimating a
population variance from a sample, and for $d=768$ the difference is 0.13\% anyway; the real
reason is compatibility, since GPT-2 was trained in TensorFlow, whose default is the biased
estimator, and we will load its weights in \cref{ch:pretrain}. Second, a quick test confirms the
contract: after normalization (with $\gamma=1,\beta=0$) every token has mean $\approx 0$ (on the
order of $10^{-8}$ in fp32, i.e., rounding noise) and variance $\approx 1$.

\begin{lstlisting}[caption={Checking the normalization contract.},label={lst:gpt:lncheck}]
torch.manual_seed(0)
x = torch.randn(2, 5, 16) * 7 + 3          # badly scaled activations
y = LayerNorm(16)(x)
print(y.mean(-1).abs().max())              # ~1e-7
print(y.var(-1, unbiased=False).mean())    # ~1.0
\end{lstlisting}

\begin{systemslens}[LayerNorm is a bandwidth problem, not a FLOP problem]
\Cref{eq:gpt:ln} costs roughly $8$ FLOPs per element but must read every element of
$\vect{x}$ and write every element of the output. Its arithmetic intensity is therefore on the
order of 1--2 FLOP/byte, two orders of magnitude below the ridge point of any modern GPU
(\cref{ch:pytorch}): the kernel runs at memory bandwidth, and the only way to make it faster is to
move fewer bytes. Our explicit implementation is about as bad as it can be---\code{mean},
\code{var}, the subtraction, the division, the scale and the shift each launch a separate kernel
that makes a full round trip through HBM. A fused LayerNorm kernel reads $\vect{x}$ once (computing
mean and variance in a single pass with Welford's algorithm or the $E[x^2]-E[x]^2$ identity in
fp32), keeps the row in registers or shared memory, and writes once. Even better is to fuse the
norm into the \emph{neighbouring} matmul or residual add so the normalized tensor never touches
HBM at all; this is one of the main wins of \code{torch.compile} (\cref{ch:pretrain}) and of the
Triton kernels in \cref{app:triton}.
\end{systemslens}

\begin{napkinmath}[How long should one LayerNorm take on an H100?]
For a $B\times T\times\dmodel = 8\times1024\times768$ bf16 activation, a fused kernel reads
$\SI{12.6}{\mega\byte}$ and writes $\SI{12.6}{\mega\byte}$. At \SI{3.35}{\tera\byte\per\second}
that is $\SI{25.2}{\mega\byte}/\SI{3.35}{\tera\byte\per\second}\approx\SI{7.5}{\micro\second}$.
The unfused version with six kernels moves roughly six times as many bytes and pays six kernel
launches of a few microseconds each: easily \SIrange{40}{60}{\micro\second}. Multiply by
$2L+1=25$ norms per forward pass and the difference is a millisecond per step---comparable to
the time of the attention matmuls at this size.
\end{napkinmath}

\begin{deepdive}[Why not batch normalization?]
Batch norm normalizes each feature using statistics pooled over the batch. For language models
this is a poor fit for three systems reasons: (i) at inference time the batch is arbitrary
(often a single request), so batch norm must fall back to running averages collected during
training, and train/test behaviour diverges; (ii) sequences have variable length, and padded
positions would pollute the statistics; (iii) under data parallelism, the statistics either
differ per GPU or require an extra all-reduce in every layer (``SyncBatchNorm''). Layer norm has
none of these problems because it never looks across tokens. Its cousin RMSNorm, which drops the
mean subtraction and the shift, is the default in modern LLMs (\cref{sec:modern:rmsnorm}).
\end{deepdive}

% =====================================================================
\section{A feed-forward network with GELU activations}\label{sec:gpt:ffn}

Attention mixes information \emph{across} positions; the feed-forward network (FFN, also called
the MLP) transforms each position \emph{independently}. It is a two-layer perceptron that first
expands the $\dmodel$-dimensional token vector to $\dff=4\dmodel$ dimensions, applies a
nonlinearity, and projects back:
\begin{equation}
\mathrm{FFN}(\vect x) = \mat W_2\,\GELU(\mat W_1\vect x+\vect b_1)+\vect b_2,\qquad
\mat W_1\in\R^{4d\times d},\ \mat W_2\in\R^{d\times 4d}.
\label{eq:gpt:ffn}
\end{equation}
The expansion gives the network a wide hidden space in which features can be computed and
gated; the contraction brings the result back to the width of the residual stream so blocks can be
stacked. Mechanistic-interpretability work often describes the FFN as a key--value memory: rows of
$\mat W_1$ act as pattern detectors (``keys'') and the corresponding columns of $\mat W_2$ as the
vectors written into the residual stream when the pattern fires (``values'').

\paragraph{GELU.}
GPT-2 uses the Gaussian error linear unit (Hendrycks and Gimpel, 2016) instead of ReLU.
$\GELU(x)=x\,\Phi(x)$, where $\Phi$ is the standard normal CDF: the input is weighted by the
probability that a standard Gaussian falls below it. GPT-2 was trained with the cheaper tanh
approximation
\begin{equation}
\GELU(x)\approx \tfrac12 x\left(1+\tanh\!\left[\sqrt{2/\pi}\,\big(x+0.044715\,x^3\big)\right]\right),
\label{eq:gpt:gelu}
\end{equation}
and since we will load its weights we must use the same formula. \Cref{fig:gpt:gelu} compares the
shapes. Unlike ReLU, GELU is smooth at the origin and lets small negative inputs pass a small
negative output (its minimum is about $-0.17$ near $x\approx-0.75$), so neurons with negative
pre-activations still receive gradient. Those properties tend to make optimization of deep
transformers slightly easier; the SiLU/Swish curve, used in the gated FFNs of modern models
(\cref{sec:modern:swiglu}), is almost indistinguishable.

\begin{figure}[t]
\centering
\begin{tikzpicture}
\begin{axis}[width=0.62\linewidth,height=5.2cm,xmin=-4,xmax=3,ymin=-0.6,ymax=3,
  grid=both,grid style={gray!20},xlabel={$x$},ylabel={activation},
  legend pos=north west,legend style={font=\scriptsize,draw=none,fill=white,fill opacity=0.8},
  tick label style={font=\scriptsize},label style={font=\small}]
\addplot[warnred,thick,domain=-4:3,samples=2] {max(0,x)};
\addlegendentry{ReLU}
\addplot[sysblue,very thick,domain=-4:3,samples=120]
  {0.5*x*(1+tanh(sqrt(2/pi)*(x+0.044715*x^3)))};
\addlegendentry{GELU (tanh approx.)}
\addplot[keygreen,thick,dashed,domain=-4:3,samples=120] {x/(1+exp(-x))};
\addlegendentry{SiLU}
\end{axis}
\end{tikzpicture}
\caption{ReLU, GELU and SiLU. GELU and SiLU are smooth, non-monotonic, and give small negative
outputs for moderately negative inputs. (Plotted analytically.)}
\label{fig:gpt:gelu}
\end{figure}

\begin{lstlisting}[caption={GELU and the expand--contract feed-forward network.},label={lst:gpt:ffn}]
import math

class GELU(nn.Module):
    def forward(self, x):
        c = math.sqrt(2.0 / math.pi)
        return 0.5 * x * (1.0 + torch.tanh(c * (x + 0.044715 * x.pow(3))))

class FeedForward(nn.Module):
    def __init__(self, emb_dim):
        super().__init__()
        self.layers = nn.Sequential(
            nn.Linear(emb_dim, 4 * emb_dim),   # expand:   (B,T,d)  -> (B,T,4d)
            GELU(),
            nn.Linear(4 * emb_dim, emb_dim),   # contract: (B,T,4d) -> (B,T,d)
        )

    def forward(self, x):
        return self.layers(x)
\end{lstlisting}

\begin{figure}[t]
\centering
\begin{tikzpicture}[x=1cm,y=1cm]
\tikzset{vec/.style={draw,fill=#1,minimum height=0.36cm,inner sep=0pt}}
\node[vec=fdata!25,minimum width=1.6cm] (x) at (0,0) {};
\node[font=\scriptsize,below=1pt of x] {$\vect x\in\R^{768}$};
\node[vec=fcompute!25,minimum width=6.4cm] (h) at (5.3,0) {};
\node[font=\scriptsize,below=1pt of h] {hidden $\in\R^{3072}$};
\node[vec=fcompute!45,minimum width=6.4cm] (g) at (5.3,1.2) {};
\node[font=\scriptsize,above=1pt of g] {$\GELU(\cdot)$ elementwise};
\node[vec=fdata!25,minimum width=1.6cm] (y) at (10.6,1.2) {};
\node[font=\scriptsize,above=1pt of y] {$\vect y\in\R^{768}$};
\draw[-{Stealth},thick] (x.east)--node[above,font=\scriptsize]{$\mat W_1$: $768\!\times\!3072$}(h.west);
\draw[-{Stealth},thick] (h.north)--(g.south);
\draw[-{Stealth},thick] (g.east)--node[above,font=\scriptsize,pos=0.4]{$\mat W_2$}(y.west);
\end{tikzpicture}
\caption{The position-wise FFN expands each token vector fourfold, applies GELU, and contracts it
back, so it can be stacked like any other shape-preserving layer.}
\label{fig:gpt:ffn}
\end{figure}

\begin{systemslens}[Where the FLOPs live: the FFN's GEMMs]
Per block, the two FFN matrices hold $8\dmodel^2$ weights versus $4\dmodel^2$ in attention's
four projections, so the FFN owns two thirds of every block's matmul FLOPs. Applied to $N=BT$
tokens at once, the first layer is a single GEMM $(N\times d)\cdot(d\times 4d)$ costing $8Nd^2$
FLOPs and moving (in bf16) $2Nd+8d^2+8Nd$ bytes (read activations, read weights, write output).
Its arithmetic intensity
\[
I(N)=\frac{8Nd^2}{10Nd+8d^2}\ \xrightarrow{N\to\infty}\ 0.8\,d
\]
grows with the number of tokens in the GEMM and saturates at $0.8d$ FLOP/byte. Whether this
kernel is compute- or memory-bound is therefore decided by \emph{how many tokens you batch
together}, a fact that will reappear as the central tension of LLM inference (\cref{ch:inference}):
during training or prefill $N$ is in the thousands and the FFN is compute bound; during
single-sequence decoding $N=1$, the intensity collapses to about 1 FLOP/byte, and the GEMM becomes
a pure weight-streaming operation. \Cref{fig:gpt:intensity} plots $I(N)$ against the H100 ridge
point.
\end{systemslens}

\begin{figure}[t]
\centering
\begin{tikzpicture}
\begin{axis}[width=0.75\linewidth,height=5.6cm,xmode=log,ymode=log,log basis x=2,
  xmin=1,xmax=16384,ymin=0.5,ymax=5000,xlabel={tokens per GEMM $N=BT$},
  ylabel={FLOP / byte},grid=major,grid style={gray!20},
  legend pos=north west,legend style={font=\scriptsize,draw=none,fill=white,fill opacity=0.85},
  tick label style={font=\scriptsize},label style={font=\small}]
\addplot[sysblue,very thick,domain=1:16384,samples=100] {8*x*768^2/(10*x*768+8*768^2)};
\addlegendentry{$\dmodel=768$ (GPT-2 small)}
\addplot[keygreen,very thick,domain=1:16384,samples=100] {8*x*4096^2/(10*x*4096+8*4096^2)};
\addlegendentry{$\dmodel=4096$}
\addplot[warnred,thick,dashed,domain=1:16384,samples=2] {295};
\addlegendentry{H100 ridge $\approx 295$}
\end{axis}
\end{tikzpicture}
\caption{Arithmetic intensity of the first FFN GEMM (bf16) as a function of how many tokens it
processes at once (analytic). Below the dashed ridge line the kernel is memory bound. The curve for
$\dmodel=768$ crosses the ridge at about 570 tokens and for $\dmodel=4096$ at about 320.}
\label{fig:gpt:intensity}
\end{figure}

\begin{napkinmath}[When does the FFN become compute bound?]
Setting $I(N)$ equal to the H100's ridge point $989/3.35\approx295$ FLOP/byte and solving
$8Nd^2=295\,(10Nd+8d^2)$ for $N$ gives $N=2360d/(8d-2950)$: about \textbf{567 tokens} for $d=768$ and
\textbf{324 tokens} for $d=4096$. A training micro-batch of $8\times1024$ tokens is far beyond
both; a decode step with batch 32 is far below. Wider models reach the ridge with fewer tokens
because the weights are re-used more per byte loaded. Note also that GPT-2 small can \emph{never}
exceed $0.8\times768\approx614$ FLOP/byte in this GEMM, only twice the ridge point: small models
are never very far from the memory wall.
\end{napkinmath}

% =====================================================================
\section{Adding shortcut connections}\label{sec:gpt:shortcut}

A deep stack of layers computes a composition $f_L\circ\cdots\circ f_1$, and by the chain rule
the gradient reaching layer 1 is a product of $L$ Jacobians. If those Jacobians typically shrink
vectors (saturating sigmoids are the textbook case), the product decays exponentially and the
early layers stop learning: the \emph{vanishing gradient} problem. Shortcut (residual, skip)
connections, introduced for very deep image networks by He et al.~(2016), replace each layer
$x\mapsto f(x)$ by $x\mapsto x+f(x)$. The Jacobian of a residual layer is $\mat I+\partial f/\partial x$,
so the backward pass always contains an identity path from the loss straight to every layer.

We can watch this happen. \Cref{lst:gpt:vanish} builds a 16-layer MLP of width 32 with sigmoid
activations, once plain and once with residual connections, runs one backward pass on random data,
and prints the mean absolute gradient of each layer's weights.

\begin{lstlisting}[caption={A vanishing-gradient experiment (our own setup).},label={lst:gpt:vanish}]
class DeepMLP(nn.Module):
    def __init__(self, depth=16, width=32, residual=False):
        super().__init__()
        self.residual = residual
        self.layers = nn.ModuleList(
            nn.Sequential(nn.Linear(width, width), nn.Sigmoid())
            for _ in range(depth))
        self.head = nn.Linear(width, 1)

    def forward(self, x):
        for layer in self.layers:
            x = x + layer(x) if self.residual else layer(x)
        return self.head(x)

x, y = torch.randn(64, 32), torch.randn(64, 1)
for residual in (False, True):
    torch.manual_seed(1)
    net = DeepMLP(residual=residual)
    nn.functional.mse_loss(net(x), y).backward()
    grads = [l[0].weight.grad.abs().mean().item() for l in net.layers]
    print(residual, [f"{g:.1e}" for g in grads[::4]])
\end{lstlisting}

\begin{figure}[t]
\centering
\begin{tikzpicture}
\begin{axis}[width=0.7\linewidth,height=5.2cm,ymode=log,xmin=0.5,xmax=16.5,ymin=1e-17,ymax=1,
  xlabel={layer index (1 = closest to input)},ylabel={mean $|\partial\mathcal L/\partial W|$},
  xtick={1,4,8,12,16},grid=major,grid style={gray!20},legend pos=south east,
  legend style={font=\scriptsize,draw=none},tick label style={font=\scriptsize},label style={font=\small}]
\addplot[warnred,very thick,mark=*] coordinates {(1,2.8e-16) (4,2.3e-14) (8,9.7e-11) (12,2.8e-7) (16,8.9e-4)};
\addlegendentry{plain stack}
\addplot[keygreen,very thick,mark=square*] coordinates {(1,3.9e-3) (4,1.9e-2) (8,3.5e-2) (12,4.1e-2) (16,4.1e-2)};
\addlegendentry{with shortcuts}
\end{axis}
\end{tikzpicture}
\caption{Gradient magnitude per layer in a 16-layer sigmoid MLP, measured on CPU in this
environment with \cref{lst:gpt:vanish}. Without shortcuts the gradient shrinks by about 12 orders
of magnitude between the last and first layer; with them it stays within one order of magnitude.}
\label{fig:gpt:vanish}
\end{figure}

\Cref{fig:gpt:vanish} shows the result: in the plain network the first layer receives gradients
around $10^{-16}$, which in fp32 arithmetic is effectively zero, while in the residual network every
layer gets a usable signal.

\paragraph{The residual stream.}
In a transformer, residual connections do more than fix gradients; they define the architecture's
central data structure. It is helpful to picture a single vector per token---the residual
stream---flowing from the embedding to the output head, with every attention and FFN sublayer
\emph{reading} from it (through a LayerNorm) and \emph{adding} a write back into it
(\cref{fig:gpt:residualstream}). Because sublayers only add, information placed in the stream
early (e.g., the token identity) remains available to all later layers unless something actively
cancels it.

\begin{figure}[t]
\centering
\resizebox{\linewidth}{!}{%
\begin{tikzpicture}[
  sub/.style={draw,rounded corners=2pt,minimum width=1.5cm,minimum height=0.6cm,font=\scriptsize,align=center},
  plus/.style={circle,draw,inner sep=0.5pt,font=\scriptsize}]
\draw[line width=3pt,fdata!60] (0,0)--(13.3,0);
\node[font=\scriptsize,anchor=east] at (0,0) {embedding};
\node[font=\scriptsize,anchor=west] at (13.3,0) {final LN, head};
\foreach \i/\x/\lab in {1/1.5/{LN\\attn}, 2/4.5/{LN\\FFN}, 3/7.5/{LN\\attn}, 4/10.5/{LN\\FFN}}{
  \node[sub,fill=fcompute!20] (n\i) at (\x+0.75,1.3) {\lab};
  \node[plus,fill=white] (p\i) at (\x+1.9,0) {$+$};
  \draw[-{Stealth},thick] (\x,0)|-(n\i.west);
  \draw[-{Stealth},thick] (n\i.east)-|(p\i.north);
}
\node[font=\scriptsize,text=srcgray] at (6.2,-0.55) {residual stream: each sublayer reads (via LN) and adds an update};
\end{tikzpicture}}
\caption{The residual-stream view of a transformer: two blocks shown, each with an attention and
an FFN sublayer reading the stream and adding their output back.}
\label{fig:gpt:residualstream}
\end{figure}

\begin{deepdive}[Pre-LN, post-LN, and residual scaling]
The original transformer (Vaswani et al., 2017) placed the normalization \emph{after} the addition,
$x\leftarrow\LN(x+f(x))$ (``post-LN''). GPT-2 moved it \emph{before} the sublayer,
$x\leftarrow x+f(\LN(x))$ (``pre-LN''). Xiong et al.~(2020) showed that post-LN produces large
gradients near the output at initialization and needs a careful learning-rate warmup, whereas
pre-LN keeps the identity path clean (nothing ever normalizes the stream itself) and trains stably
at depth. The price is that the stream's magnitude grows with depth: each of the $2L$ sublayers adds
a roughly independent update, so the variance of the stream grows linearly in $L$. GPT-2 therefore
scales the initial weights of each sublayer's \emph{output} projection by $1/\sqrt{2L}$, and adds
a final LayerNorm before the head so the logits see a normalized vector. Both conventions survive
in today's models.
\end{deepdive}

% =====================================================================
\section{Connecting attention and linear layers in a transformer block}\label{sec:gpt:block}

We now have every ingredient of the GPT-2 block (\cref{fig:gpt:block}): masked multi-head
attention from \cref{ch:attention}, the FFN, two LayerNorms, two residual connections, and dropout
applied to each sublayer's output before it is added back (sometimes called \emph{residual
dropout}). The block is shape-preserving: $(B,T,\dmodel)$ in, $(B,T,\dmodel)$ out. That property is
not incidental---it is what allows the block to be stacked any number of times and what makes the
model's depth a single configuration number.

\begin{figure}[t]
\centering
\begin{tikzpicture}[
  op/.style={draw,rounded corners=2pt,minimum width=4.0cm,minimum height=0.6cm,font=\small},
  plus/.style={circle,draw,fill=white,inner sep=0.6pt,font=\small},
  arr/.style={-{Stealth[length=2mm]},thick}]
\node[font=\small] (in) at (0,0) {$x$ \ (B, T, d)};
\node[op,fill=white] (ln1) at (0,1.1) {LayerNorm 1};
\node[op,fill=fcompute!25] (mha) at (0,2.0) {Masked multi-head attention};
\node[op,fill=gray!10] (d1) at (0,2.9) {Dropout};
\node[plus] (a1) at (0,3.75) {$+$};
\node[op,fill=white] (ln2) at (0,4.6) {LayerNorm 2};
\node[op,fill=fcompute!25] (ff) at (0,5.5) {Feed-forward (GELU)};
\node[op,fill=gray!10] (d2) at (0,6.4) {Dropout};
\node[plus] (a2) at (0,7.25) {$+$};
\node[font=\small] (out) at (0,8.15) {$x'$ \ (B, T, d)};
\draw[arr] (in)--(ln1); \draw[arr] (ln1)--(mha); \draw[arr] (mha)--(d1); \draw[arr] (d1)--(a1);
\draw[arr] (a1)--(ln2); \draw[arr] (ln2)--(ff); \draw[arr] (ff)--(d2); \draw[arr] (d2)--(a2); \draw[arr] (a2)--(out);
\draw[arr,keygreen] (0,0.45) -- (-2.6,0.45) |- (a1.west);
\draw[arr,keygreen] (0,4.15) -- (-2.6,4.15) |- (a2.west);
\node[font=\scriptsize,text=keygreen,rotate=90] at (-2.85,2.1) {shortcut};
\node[font=\scriptsize,text=keygreen,rotate=90] at (-2.85,5.7) {shortcut};
\node[font=\scriptsize,align=left,anchor=west,text=srcgray] at (2.3,2.0) {mixes information\\across positions};
\node[font=\scriptsize,align=left,anchor=west,text=srcgray] at (2.3,5.5) {transforms each\\position independently};
\end{tikzpicture}
\caption{The pre-LN GPT-2 transformer block. Green arrows are the two shortcut connections.}
\label{fig:gpt:block}
\end{figure}

\begin{lstlisting}[caption={The transformer block (cf.\ \ident{llmsys/gpt.py}).},label={lst:gpt:block}]
from llmsys.layers import MultiHeadAttention

class TransformerBlock(nn.Module):
    def __init__(self, cfg):
        super().__init__()
        self.att = MultiHeadAttention(
            d_in=cfg["emb_dim"], d_out=cfg["emb_dim"],
            context_length=cfg["context_length"], num_heads=cfg["n_heads"],
            dropout=cfg["drop_rate"], qkv_bias=cfg["qkv_bias"])
        self.ff = FeedForward(cfg["emb_dim"])
        self.norm1 = LayerNorm(cfg["emb_dim"])
        self.norm2 = LayerNorm(cfg["emb_dim"])
        self.drop_shortcut = nn.Dropout(cfg["drop_rate"])

    def forward(self, x):
        x = x + self.drop_shortcut(self.att(self.norm1(x)))   # attention
        x = x + self.drop_shortcut(self.ff(self.norm2(x)))    # FFN sublayer
        return x
\end{lstlisting}

The code mirrors the figure line by line. Note that we never write \code{shortcut = x} and add it
later: \code{x = x + f(x)} creates a \emph{new} tensor, so the old \code{x} is still referenced by
the autograd graph and the in-place semantics are safe.

\begin{systemslens}[Anatomy of one block: what is compute bound and what is not]
\Cref{tab:gpt:blockops} lists the operations in one GPT-2 small block for a training micro-batch of
$N=BT=8192$ tokens in bf16, with their FLOPs and the minimum bytes they must move. Four GEMM groups
(QKV, attention scores and weighted sum, output projection, two FFN layers) are comfortably compute
bound; everything else---two LayerNorms, softmax, GELU, two dropouts, two residual adds---is
elementwise or a row reduction with intensity of a few FLOP/byte. In an eager PyTorch run, those
cheap operations are a long tail of small memory-bound kernels, and on a GPU they can take 30--50\% of
the step time despite contributing under 1\% of the FLOPs. This is the motivation for fusion:
FlashAttention fuses mask+softmax+dropout into the attention matmuls (\cref{ch:attention}),
GEMM epilogues fuse bias+GELU, and compilers fuse residual-add+LayerNorm.
\end{systemslens}

\begin{table}[t]
\centering\small
\begin{tabular}{lrrl}
\toprule
Operation (one block) & FLOPs & Min.\ bytes moved & Bound\\
\midrule
Q, K, V projections ($3\times d\!\times\! d$) & $6Nd^2$ & $6Nd+6d^2+\ldots$ & compute\\
$QK^\top$ and $AV$ (per sequence $T\times T$) & $4NTd$ & $\propto BHT^2$ if unfused & compute\\
mask $+$ softmax $+$ dropout & $\approx 5BHT^2$ & $\approx 4BHT^2\cdot 2$ & memory\\
output projection & $2Nd^2$ & $4Nd+2d^2$ & compute\\
FFN up $+$ down & $16Nd^2$ & $20Nd+16d^2$ & compute\\
GELU (tanh) & $\approx 15\cdot 4Nd$ & $16Nd$ & memory\\
2$\times$ LayerNorm, 2$\times$ residual add, 2$\times$ dropout & $\approx 30Nd$ & $\approx 24Nd$ & memory\\
\bottomrule
\end{tabular}
\caption{Operation inventory for one pre-LN block on $N=BT$ tokens ($d=\dmodel$, bf16, 2 bytes per
element). Counts are approximate for the elementwise rows.}
\label{tab:gpt:blockops}
\end{table}

% =====================================================================
\section{Coding the GPT model}\label{sec:gpt:model}

With the block in hand, the full model is short. \Cref{lst:gpt:model} replaces the placeholders of
the skeleton with real blocks and a real final LayerNorm. (The library version in
\ident{llmsys/gpt.py} additionally threads an optional KV cache and padding mask through the
blocks; we will need those from \cref{ch:inference} onwards.)

\begin{lstlisting}[caption={The complete GPT model.},label={lst:gpt:model}]
class GPTModel(nn.Module):
    def __init__(self, cfg):
        super().__init__()
        self.cfg = cfg
        self.tok_emb = nn.Embedding(cfg["vocab_size"], cfg["emb_dim"])
        self.pos_emb = nn.Embedding(cfg["context_length"], cfg["emb_dim"])
        self.drop_emb = nn.Dropout(cfg["drop_rate"])
        self.trf_blocks = nn.ModuleList(
            [TransformerBlock(cfg) for _ in range(cfg["n_layers"])])
        self.final_norm = LayerNorm(cfg["emb_dim"])
        self.out_head = nn.Linear(cfg["emb_dim"], cfg["vocab_size"], bias=False)

    def forward(self, idx):                        # (B, T) token ids
        B, T = idx.shape
        pos = torch.arange(T, device=idx.device)
        x = self.drop_emb(self.tok_emb(idx) + self.pos_emb(pos))
        for blk in self.trf_blocks:
            x = blk(x)
        return self.out_head(self.final_norm(x))   # (B, T, V) logits
\end{lstlisting}

\paragraph{Counting parameters.}
Instantiating the model with the 124M configuration and summing \code{p.numel()} over
\code{model.parameters()} gives \num{163009536}, not 124 million. The difference is
\emph{weight tying}: the released GPT-2 reuses the token-embedding matrix ($\vocab\times\dmodel$) as
the output projection, while our model has a separate \code{out\_head} of the same shape.
Removing it subtracts $\num{50257}\times768=\num{38597376}$, leaving \num{124412160}; the
remaining \num{27648} difference to the official \num{124439808} comes from the QKV biases the
released checkpoints use ($3\times768\times12$). Every parameter can be accounted for in closed
form. Per block,
\begin{align}
P_\text{attn} &= \underbrace{3d^2}_{W_Q,W_K,W_V} + \underbrace{3d}_{\text{QKV bias, optional}} + \underbrace{d^2+d}_{\text{out proj}}, \\
P_\text{FFN}  &= \underbrace{4d^2+4d}_{W_1,b_1} + \underbrace{4d^2+d}_{W_2,b_2},\qquad
P_\text{LN} = 2\times 2d,
\end{align}
and for the whole model
\begin{equation}
P = \underbrace{\vocab d}_{\text{tok emb}} + \underbrace{T_\text{max} d}_{\text{pos emb}}
 + L\,(P_\text{attn}+P_\text{FFN}+P_\text{LN}) + \underbrace{2d}_{\text{final LN}}
 + \underbrace{\vocab d}_{\text{head, if untied}} \;\approx\; 12Ld^2 + 2\vocab d .
\label{eq:gpt:params}
\end{equation}
Lab~5 asks you to implement this formula and asserts that it matches the module exactly for
several configurations. \Cref{fig:gpt:parambreakdown} shows how the budget splits across the
four GPT-2 sizes: the embedding and head dominate the small model, while for GPT-2 XL the
transformer blocks account for 90\% of the parameters and the FFN alone for 60\%.

\begin{figure}[t]
\centering
\begin{tikzpicture}
\begin{axis}[ybar stacked,width=0.62\linewidth,height=5.6cm,bar width=22pt,
  symbolic x coords={small,medium,large,XL},xtick=data,ymin=0,ymax=1750,
  ylabel={parameters (millions)},legend style={font=\scriptsize,at={(1.03,0.5)},anchor=west,draw=none},reverse legend,
  tick label style={font=\scriptsize},label style={font=\small},enlarge x limits=0.18]
\addplot[fill=fdata!55,draw=fdata] coordinates {(small,39.4) (medium,52.5) (large,65.6) (XL,82.0)};
\addplot[fill=fcompute!35,draw=fcompute] coordinates {(small,28.3) (medium,100.7) (large,236.0) (XL,491.6)};
\addplot[fill=fcompute!75,draw=fcompute] coordinates {(small,56.7) (medium,201.4) (large,472.1) (XL,983.4)};
\addplot[fill=napkin!45,draw=napkin] coordinates {(small,38.6) (medium,51.5) (large,64.3) (XL,80.4)};
\legend{tok $+$ pos embeddings, attention, FFN, untied output head}
\end{axis}
\end{tikzpicture}
\caption{Parameter breakdown of the four GPT-2 sizes in our untied configuration (LayerNorm
parameters are too small to see). Computed with \code{count\_params\_analytic} from Lab~5.}
\label{fig:gpt:parambreakdown}
\end{figure}

\begin{napkinmath}[Memory footprint of the weights]
At 4 bytes per fp32 parameter, the untied 163M model needs
$\num{163009536}\times4\,\text{B}\approx\SI{622}{\mebi\byte}$; the tied 124M checkpoint in bf16 needs
$\SI{249}{\mega\byte}$. GPT-2 XL in fp32 is \SI{6.1}{\gibi\byte}. These numbers are only the
\emph{weights}: training with AdamW in mixed precision needs about 16 bytes per parameter for
weights, gradients and optimizer state (\cref{ch:pretrain}), so GPT-2 XL needs
$\approx\SI{26}{\giga\byte}$ before a single activation is stored---already more than a 24\,GB
consumer GPU.
\end{napkinmath}

\begin{systemslens}[Weight tying as a systems trade-off]
Tying saves $\vocab\dmodel$ parameters (24\% of GPT-2 small, 5\% of GPT-2 XL) and the
corresponding gradient and optimizer memory. It does \emph{not} save FLOPs: the output projection
is still a $d\times\vocab$ GEMM on every position. It also couples two very different gradient
signals (sparse row updates from the embedding lookup, dense updates from the head) into one
tensor, and in pipeline-parallel training the first and last stages must synchronize that shared
tensor's gradient with an extra all-reduce. Small models (GPT-2, Qwen3-0.6B to 4B) tend to tie;
large ones (Qwen3-8B and up, Llama 3) usually do not.
\end{systemslens}

\subsection{Initialization}
A freshly constructed \code{nn.Linear} uses PyTorch's default Kaiming-uniform initialization, which
was designed for ReLU convolutional networks. GPT-2 instead draws every weight matrix from
$\mathcal N(0,0.02^2)$, sets biases to zero, and (as discussed in \cref{sec:gpt:shortcut}) shrinks the
standard deviation of the two \emph{output} projections in each block (attention's \code{out\_proj}
and the FFN's second linear) by $1/\sqrt{2L}$, so that the $2L$ residual updates sum to a stream of
roughly constant scale at initialization. \Cref{lst:gpt:init} applies this recipe. It matters only when
training from scratch---loaded checkpoints overwrite every value---but a bad initialization is one of
the most common reasons a from-scratch pretraining run diverges in its first few hundred steps.

\begin{lstlisting}[caption={GPT-2 style initialization.},label={lst:gpt:init}]
import math

def init_gpt2_weights(model, std=0.02):
    L = model.cfg["n_layers"]
    for name, p in model.named_parameters():
        if p.dim() == 2:                          # weight matrices and embeddings
            s = std / math.sqrt(2 * L) if name.endswith(
                ("out_proj.weight", "ff.layers.2.weight")) else std
            torch.nn.init.normal_(p, mean=0.0, std=s)
        elif name.endswith(("bias", "shift")):
            torch.nn.init.zeros_(p)
        # LayerNorm scales keep their default value of 1
\end{lstlisting}

\begin{systemslens}[Initialization at scale]
For a 124M-parameter model initialization takes a fraction of a second; for a 70B-parameter model,
drawing 70 billion Gaussian samples on one CPU core takes minutes and needs 280\,GB of host memory in
fp32. Large-scale training frameworks therefore construct the model on PyTorch's \code{meta} device
(shapes only, no storage), shard it across GPUs first (FSDP, \cref{ch:distributed}), and only then run
the initializer on each GPU's shard. Two subtleties follow. First, every data-parallel replica must
start from \emph{identical} weights, so either one rank initializes and broadcasts, or all ranks seed
their random generators identically. Second, the initializer must be written per-parameter-name (as in
\cref{lst:gpt:init}) rather than per-module-constructor, because under sharding the module that
``owns'' a weight may never see its full shape.
\end{systemslens}

\begin{deepdive}[What dropout costs]
Dropout looks free, but in an eager implementation each dropout call generates a random tensor of the
same shape as its input, stores a boolean mask for the backward pass (1 byte per element, or one bit if
packed), and launches two extra kernels. GPT-2 has $2L+1$ residual/embedding dropouts plus $L$ dropouts
on the $BHT^2$ attention probabilities, the largest activation tensor in the model. Fused kernels
regenerate the mask in the backward pass from the random generator's seed and offset (Philox counters)
instead of storing it, which is how FlashAttention supports attention dropout without materializing
anything. For pretraining modern LLMs on trillions of tokens, the usual answer is simpler: set the
dropout rate to zero (\cref{ch:modern}).
\end{deepdive}

% =====================================================================
\section{Generating text}\label{sec:gpt:generate}

The model outputs logits for every position, but to generate text we only need the last one: it
scores every candidate for the token that follows the whole input. Generation is a loop
(\cref{fig:gpt:genloop}): run the model, take the logits at position $-1$, pick a token, append it
to the input, repeat. Picking the token with the highest score is \emph{greedy decoding}. Because
softmax is monotonic, the argmax of the logits equals the argmax of the probabilities, so the softmax
can be skipped; we will need the probabilities once we sample (\cref{ch:pretrain,ch:scaling}).

\begin{figure}[t]
\centering
\begin{tikzpicture}[
  box/.style={draw,rounded corners=2pt,minimum height=0.65cm,font=\small,align=center},
  arr/.style={-{Stealth[length=2mm]},thick}]
\node[box,fill=fdata!20,minimum width=3.4cm] (ctx) at (0,0) {Hello , I am};
\node[box,fill=fcompute!25,minimum width=1.8cm] (gpt) at (3.6,0) {GPT};
\node[box,fill=white,minimum width=2.6cm] (last) at (7.0,0) {logits at pos.\ $-1$\\ \scriptsize $(B,\vocab)$};
\node[box,fill=keygreen!20,minimum width=1.7cm] (am) at (10.1,0) {argmax};
\node[box,fill=fdata!35,minimum width=1.4cm] (tok) at (10.1,-1.6) {``a''};
\draw[arr] (ctx)--(gpt); \draw[arr] (gpt)--(last); \draw[arr] (last)--(am); \draw[arr] (am)--(tok);
\draw[arr,sysblue] (tok.west) -- (0,-1.6) node[midway,below,font=\scriptsize]{append; crop to the last \code{context\_size} tokens} -- (ctx.south);
\end{tikzpicture}
\caption{The autoregressive generation loop. Each iteration re-runs the model on the (cropped)
sequence, keeps only the last position's logits, and appends the chosen token.}
\label{fig:gpt:genloop}
\end{figure}

\begin{lstlisting}[caption={Greedy generation without a cache.},label={lst:gpt:generate}]
@torch.no_grad()
def generate_greedy(model, idx, max_new_tokens, context_size):
    model.eval()                               # disable dropout
    for _ in range(max_new_tokens):
        idx_cond = idx[:, -context_size:]      # never exceed the position table
        logits = model(idx_cond)[:, -1, :]     # (B, V): only the last position
        next_id = torch.argmax(logits, dim=-1, keepdim=True)
        idx = torch.cat([idx, next_id], dim=1)
    return idx
\end{lstlisting}

Three details are easy to get wrong. (i) \code{model.eval()} switches dropout off; leaving it on
makes generation nondeterministic and slightly worse. (ii) The crop
\code{idx[:, -context\_size:]} is mandatory for a model with learned absolute positions: GPT-2 has
exactly 1024 position embeddings and indexing the 1025th raises an error. (iii) Since the weights
are random, the untrained model produces fluent-looking nonsense: a valid sequence of token IDs that
decodes to unrelated word pieces. Giving the model useful weights is the subject of
\cref{ch:pretrain}, where we both train it and load OpenAI's checkpoint.

\begin{napkinmath}[The quadratic cost of cache-less generation]
Each step of \cref{lst:gpt:generate} re-processes the whole prefix. A forward pass over a sequence of
length $t$ costs about $t\cdot(2P_\text{matmul} + 4Ldt)$ FLOPs (\cref{sec:gpt:accounting}). Generating
1000 tokens from a 10-token prompt with GPT-2 small therefore costs
$\sum_{t=10}^{1009} F(t)\approx139$ TFLOP (computed with the Lab~5 formula), whereas
processing each new token only once---which is what the KV cache of \cref{ch:inference} enables---costs
about 0.27 TFLOP, roughly \textbf{500$\times$} less. The cache-less loop is
$O(n^2)$ in the matmuls and $O(n^3)$ in attention; it is fine for a teaching example and useless in
production.
\end{napkinmath}

% =====================================================================
\section{Accounting for every FLOP and byte}\label{sec:gpt:accounting}

We close the chapter with the systems analysis that the rest of the book builds on. The goal is
three formulas---FLOPs, weight memory, activation memory---each validated against PyTorch.

\subsection{FLOPs of a forward pass}
A matrix product $(m\times k)\cdot(k\times n)$ performs $mkn$ multiply--accumulates, conventionally
counted as $2mkn$ FLOPs. Every \code{nn.Linear} applied to $N$ tokens is such a product with
$m=N$, so it costs $2N$ times its number of weights. Summing over the block gives $2N\cdot12d^2$; the
head adds $2N\vocab d$. The two attention matmuls are different because their ``weights'' are other
activations: $QK^\top$ for one head of one sequence is $(T\times\dhead)\cdot(\dhead\times T)$, i.e.
$2T^2\dhead$ FLOPs, and $AV$ costs the same, so per layer and sequence the pair costs
$4T^2\dhead H = 4T^2 d$. Putting it together for $B$ sequences of length $T$:
\begin{equation}
F_\text{fwd}(B,T) = BT\Big[\,L\,(24d^2 + 4Td) + 2\vocab d\,\Big].
\label{eq:gpt:flops}
\end{equation}
Per token, the matmul part is twice the number of non-embedding parameters (plus the head), the
classic ``$2N$ FLOPs per token'' rule; the attention part grows linearly with context and overtakes
the $24d^2$ term once $T>6d$ (4608 tokens for GPT-2 small).

\begin{lstlisting}[caption={Verifying \cref{eq:gpt:flops} against PyTorch's own FLOP counter.},label={lst:gpt:flopcounter}]
from torch.utils.flop_counter import FlopCounterMode

def forward_flops(cfg, batch, seq_len):
    d, V, L = cfg["emb_dim"], cfg["vocab_size"], cfg["n_layers"]
    per_token = L * (24 * d * d + 4 * seq_len * d) + 2 * d * V
    return batch * seq_len * per_token

model = GPTModel(cfg).eval()
idx = torch.randint(0, cfg["vocab_size"], (2, 16))
with FlopCounterMode(display=False) as fc:
    model(idx)
assert fc.get_total_flops() == forward_flops(cfg, 2, 16)   # exact match
\end{lstlisting}

The match in \cref{lst:gpt:flopcounter} is \emph{exact} because \code{FlopCounterMode} counts only
matmul-like operators with the same $2mkn$ convention, and our explicit attention computes the full
$T\times T$ score matrix even though half of it is masked. (A causal FlashAttention kernel skips the
masked tiles and does half the attention FLOPs; FLOP counters then disagree with ``useful'' FLOPs, a
recurring source of confusion when people report utilization.) For GPT-2 small at $T=1024$ the
formula gives 0.285 GFLOP per token.

\begin{keyidea}[The 6N rule]
The backward pass of a matmul computes two products of the same size as the forward one (gradient
with respect to the input and with respect to the weight), so it costs about twice the forward.
Training therefore costs $\approx 3\times F_\text{fwd}$, i.e. about $6P$ FLOPs per token for a
model with $P$ (non-embedding) parameters when attention is a minor term. Training GPT-2 small on
10 billion tokens costs $6\times\num{1.24e8}\times10^{10}\approx\num{7.4e18}$ FLOPs: about
2.1 hours on one H100 at 100\% utilization, or roughly 5 hours at a realistic 40\% MFU
(\cref{ch:pretrain}).
\end{keyidea}

\subsection{Activation memory}
For the backward pass, autograd keeps every tensor that some operation's derivative needs.
Walking through our block on $s$ tokens per sequence, batch $b$, width $d$ and $h$ heads, the
saved tensors are: the input of LN1 ($sbd$ elements); its output, shared by the Q/K/V projections
($sbd$); $Q$ and $K$ for the score matmul ($2sbd$); the softmax probabilities ($bhs^2$); $V$
($sbd$); the context fed to the output projection ($sbd$); the input of LN2 ($sbd$); the FFN input
($sbd$); the pre-GELU hidden ($4sbd$); and the post-GELU hidden fed to the second linear ($4sbd$).
In total
\begin{equation}
A_\text{layer} \approx \big(16\,sbd + b h s^2\big)\ \text{elements},
\label{eq:gpt:act}
\end{equation}
which with 2-byte activations and 1-byte dropout masks becomes the widely quoted
$sbd\,(34+5hs/d)$ bytes of Korthikanti et al.~(2022). Two consequences: activation memory is
linear in batch size and in depth, and the attention term grows \emph{quadratically} in sequence
length, which is precisely what FlashAttention removes by never storing the probabilities.

\begin{lstlisting}[caption={Measuring saved activations with autograd hooks.},label={lst:gpt:savedhooks}]
def measure_saved_activation_bytes(model, idx):
    params = {p.untyped_storage().data_ptr() for p in model.parameters()}
    seen = {}
    def pack(t):                        # called for every tensor saved for backward
        ptr = t.untyped_storage().data_ptr()
        if ptr not in params and ptr not in seen:
            seen[ptr] = t.untyped_storage().nbytes()
        return t
    with torch.autograd.graph.saved_tensors_hooks(pack, lambda t: t):
        model(idx)
    return sum(seen.values())
\end{lstlisting}

Running \cref{lst:gpt:savedhooks} on the tiny test configuration ($d=64$, 2 layers, 4 heads, one
sequence of 64 tokens, fp32) gives the numbers in \cref{tab:gpt:actmeasured}. The measured value is
about $2.1\times$ the estimate of \cref{eq:gpt:act}. The explanation is instructive: our explicit
tanh-GELU is built from half a dozen primitive operations, and autograd saves several of their
intermediates, each $4d$ wide. Replacing it with PyTorch's fused \code{nn.GELU(approximate="tanh")},
which saves only its input, removes 28\% of all activation memory; the explicit LayerNorm similarly
saves extra intermediates that a fused kernel would not. Fusion saves memory, not just time.

\begin{table}[h]
\centering\small
\begin{tabular}{lr}
\toprule
Configuration (tiny GPT, $b=1$, $s=64$, fp32) & Saved activations\\
\midrule
Estimate, \cref{eq:gpt:act} & \SI{655}{\kilo\byte}\\
Measured, explicit GELU and LayerNorm & \SI{1387}{\kilo\byte}\\
Measured, fused \code{nn.GELU}, explicit LayerNorm & \SI{994}{\kilo\byte}\\
\bottomrule
\end{tabular}
\caption{Activation memory: estimate vs.\ measurement with saved-tensor hooks (measured on CPU in
this environment). The measured value doubles exactly when the batch doubles.}
\label{tab:gpt:actmeasured}
\end{table}

\begin{napkinmath}[Activations of GPT-2 small at full context]
For $b=8$, $s=1024$, $d=768$, $h=12$, $L=12$ in bf16, \cref{eq:gpt:act} gives per layer
$16\cdot1024\cdot8\cdot768\cdot2\,\text{B}\approx\SI{201}{\mega\byte}$ for the linear terms and
$8\cdot12\cdot1024^2\cdot2\,\text{B}\approx\SI{201}{\mega\byte}$ for the attention probabilities:
\SI{0.4}{\giga\byte} per layer, \SI{4.8}{\giga\byte} for the whole model, plus the
\SI{1.65}{\giga\byte} fp32 logits. That is roughly ten times the weights. At $s=8192$ the attention
term alone would be \SI{12.9}{\giga\byte} \emph{per layer}. This is why long-context training
requires both FlashAttention (no stored probabilities) and activation checkpointing
(\cref{ch:pretrain}).
\end{napkinmath}

\subsection{Where does the time go?}
FLOPs are not time. The Lab~5 \code{bench.py} script profiles a forward pass with
\code{torch.profiler} and groups the time by operator. On a GPU you will typically find the large
GEMMs (\code{aten::addmm}, \code{aten::mm}) at the top of the table but far from 100\% of the time,
followed by a long tail of elementwise kernels (\code{aten::mul}, \code{aten::add},
\code{aten::tanh}, \code{aten::pow}) from our explicit GELU and LayerNorm, and copies from
\code{transpose().contiguous()} in attention. Dividing \cref{eq:gpt:flops} by the measured time gives
the achieved FLOP/s; dividing that by the device peak gives the \emph{model FLOPs utilization}
(MFU), the single most useful efficiency number in LLM systems and the subject of
\cref{ch:pretrain}.

\begin{systemslens}[A first look at kernel fusion and compilation]
The gap between ``FLOPs'' and ``time'' in an eager model has three sources: memory-bound
elementwise kernels that each round-trip through HBM, per-kernel launch overhead (a few microseconds
each, which dominates small models), and extra copies from layout changes. All three are attacked by
the same technique---\emph{fusion}---either by hand (fused LayerNorm, fused bias-GELU,
FlashAttention), by vendor libraries (cuBLASLt epilogues), or automatically by
\code{torch.compile}, whose TorchInductor backend traces the model, groups pointwise and reduction
operations into generated Triton kernels, and can capture the whole forward pass in a CUDA graph to
eliminate launch overhead. We will measure the effect in \cref{ch:pretrain,ch:inference} and write
such kernels ourselves in \cref{app:triton}.
\end{systemslens}

% =====================================================================
\section{Summary}
\begin{itemize}
\item A GPT model is an embedding layer, a stack of $L$ identical shape-preserving transformer
  blocks, a final LayerNorm, and a linear head producing $(B,T,\vocab)$ logits.
\item LayerNorm normalizes each token across its features; it is independent of the batch and
  therefore well suited to variable-length, distributed training. It is memory bound and an
  important fusion target.
\item The FFN expands to $4\dmodel$, applies GELU, and contracts back. It holds two thirds of each
  block's weights and FLOPs; its arithmetic intensity depends on how many tokens are processed at once.
\item Residual connections keep gradients alive in deep networks and define the residual stream;
  pre-LN placement and $1/\sqrt{2L}$ output scaling keep training stable.
\item GPT-2 small has \num{163009536} parameters with a separate head and \num{124439808} in the
  released tied checkpoint; parameters, FLOPs ($\approx 2P$ per token forward, $6P$ training) and
  activation memory ($\approx 16sbd+bhs^2$ elements per layer) can all be computed exactly and checked
  against PyTorch.
\item Greedy generation repeatedly runs the model and appends the argmax token; without a cache it
  re-processes the whole prefix each step, which is quadratically wasteful.
\end{itemize}

% =====================================================================
\section{Exercises}

\begin{exercise}[GPT-2 sizes]
Using \cref{eq:gpt:params}, compute the parameter counts of GPT-2 medium, large and XL in the
untied configuration, and separately the share held by attention and by the FFN. Check your numbers
against \cref{tab:gpt:sizes} and \cref{fig:gpt:parambreakdown}. At what $\dmodel$ does the FFN first
hold more parameters than the token embedding plus the head, for $L=12$ and $\vocab=\num{50257}$?
\end{exercise}

\begin{exercise}[Separate dropout rates]
Our configuration uses one \code{drop\_rate} for the embedding dropout, the residual dropout, and
the attention-weight dropout. Change the configuration and code so the three rates can be set
independently, and argue which of the three is cheapest to remove from a systems perspective.
\end{exercise}

\begin{exercise}[Vocabulary padding]
Pad the vocabulary of GPT-2 small to the next multiple of 128. How many parameters and how many
bytes of fp32 optimizer state (two Adam moments) does that add? What must the generation code do
to make sure the padded rows are never sampled?
\end{exercise}

\begin{exercise}[Ridge point]
Derive the crossover $N^\star$ at which the first FFN GEMM becomes compute bound on an A100
(312 TFLOP/s bf16, 2.04 TB/s) for $d=768$ and $d=4096$. Repeat for the second FFN GEMM, whose
input is $4d$ wide. Which of the two GEMMs reaches the ridge first, and why?
\end{exercise}

\begin{exercise}[Residual growth]
Initialize GPT-2 small with all weights drawn from $\mathcal N(0,0.02^2)$ and measure the RMS of the
residual stream after each block on random tokens. Then scale the weights of every
\code{out\_proj} and second FFN linear by $1/\sqrt{2L}$ and measure again. Plot both curves and explain
the shape of the first.
\end{exercise}

\begin{exercise}[Post-LN]
Implement a post-LN variant of \code{TransformerBlock}. Train a 24-layer tiny model with both
variants for a few hundred steps on the same data without learning-rate warmup. Which one diverges
first as you raise the learning rate?
\end{exercise}

\begin{exercise}[Fused activations save memory]
Using \cref{lst:gpt:savedhooks}, measure saved-activation bytes for (a) our explicit model,
(b) the model with \code{nn.GELU(approximate="tanh")}, (c) additionally with \code{nn.LayerNorm}.
Explain each reduction by listing the tensors autograd stops saving.
\end{exercise}

\begin{exercise}[Causal FLOPs]
\Cref{eq:gpt:flops} counts the full $T\times T$ attention. Derive the ``useful'' FLOPs when only the
causal half is computed, and compute the relative difference for GPT-2 small at $T=1024$ and
$T=8192$. Why does it matter when comparing MFU numbers across papers?
\end{exercise}

\begin{exercise}[Cost of generation]
Write a function that returns the total FLOPs of generating $n$ tokens from a $p$-token prompt
(a) without a cache, (b) with a cache. Plot the ratio as a function of $n$ for $p=100$. At which
$n$ does the cached version save $100\times$?
\end{exercise}

\begin{exercise}[Chunked loss]
Design (and implement) a cross-entropy computation that processes the sequence in chunks of 128
positions so that the full $(B,T,\vocab)$ logits tensor is never materialized. What is the peak
memory saving for $B=8$, $T=1024$, $\vocab=\num{50257}$ in fp32? What extra cost do you pay in the
backward pass?
\end{exercise}

% =====================================================================
\section{Lab 5: Assemble GPT-2 and account for every parameter, FLOP and byte}\label{sec:gpt:lab}

\begin{labbox}{5}{Assemble GPT-2 and account for every parameter, FLOP and byte}
\labdir{lab05_gpt}
\labsection{Goals}
Build the GPT-2 architecture from its components so that it is bit-compatible with
\ident{llmsys.gpt.GPTModel}, then turn the chapter's accounting into tested code: an exact
parameter formula, an exact forward-FLOP formula, and an activation-memory measurement.

\labsection{Setup}
\begin{lstlisting}[style=shell]
cd code && pip install -r requirements.txt
pytest -q labs/lab05_gpt            # runs the tests against your starter.py
\end{lstlisting}

\labsection{Tasks}
\begin{enumerate}
\item \textbf{Components.} Implement \code{LayerNorm.forward}, \code{gelu\_tanh} and
  \code{FeedForward.\_\_init\_\_}. Keep the attribute names (\code{scale}, \code{shift},
  \code{layers}) so state dicts are interchangeable with the library.
\item \textbf{Block and model.} Implement \code{TransformerBlock.forward} (pre-LN, residual dropout)
  and \code{GPTModel.forward}. The test loads the library model's weights into yours and requires
  identical logits.
\item \textbf{Generation.} Implement \code{generate\_greedy} with context cropping; the test generates
  past \code{context\_length} to make sure you crop.
\item \textbf{Accounting.} Implement \code{count\_params\_analytic} (with and without tying) and
  \code{forward\_flops}. The tests compare against \code{sum(p.numel())} and against
  \code{FlopCounterMode} for several shapes, and pin the GPT-2 small numbers \num{163009536} and
  \num{124439808}.
\item \textbf{Activation memory.} Implement \code{activation\_bytes\_estimate} (\cref{eq:gpt:act})
  and \code{measure\_saved\_activation\_bytes} with \code{saved\_tensors\_hooks}.
\end{enumerate}

\labsection{Measurements}
\begin{itemize}
\item Run \code{python labs/lab05\_gpt/bench.py --impl starter --seq 256 --batch 4} on your machine.
  Record forward time, achieved TFLOP/s, and the top 10 operators. Then swap in the fused GELU and
  \code{nn.LayerNorm} and re-run. Expect the elementwise tail to shrink noticeably on a GPU.
\item Plot measured activation bytes against batch size (should be a straight line through the
  origin) and against sequence length (should curve upwards: the $bhs^2$ term).
\item If you have the HF \code{gpt2} \code{model.safetensors}, pass \code{--weights} and confirm
  your model generates coherent English with greedy decoding.
\end{itemize}

\labsection{Deliverables}
Passing tests; a short table of operator times before and after fusion; the two activation plots
with one paragraph explaining their shapes.

\labsection{Stretch goals}
\begin{itemize}
\item Wrap the model in \code{torch.compile} and compare time and saved-activation bytes.
\item Implement the chunked cross-entropy of the last exercise and measure peak memory with
  \code{torch.cuda.max\_memory\_allocated}.
\end{itemize}
\end{labbox}
